% GENERATED FILE. Edit canonical lessons, then run npm run generate:books.
% canonical-source-hash: fnv1a64:343da868a5e90e93
% canonical-lessons: TE-C185-kavaru, TE-C185-stampu, TE-C185-parselu, TE-C185-tapala, TE-C185-pensilu

\chapter{Envelope, Postage stamp, Parcel, Post, Pencil}
\label{ch:te-envelope-postage-stamp-parcel-post-pencil}
\hlchaptermodality{185}

\begin{chapteropening}
\textbf{By the end of this chapter:} \emph{I can say an envelope, a postage stamp, a parcel, the post and a pencil.}

Complete the last lesson of chapter 185: I can say an envelope, a postage stamp, a parcel, the post and a pencil.
\end{chapteropening}

\section[\texorpdfstring{\te{కవరు} (kavaru)}{కవరు (kavaru)}]{\te{కవరు} (kavaru) — an envelope}

\label{lesson:TE-C185-kavaru}



\begin{quote}
Before the new one: say the Telugu for a cucumber, then the Telugu for okra.
\end{quote}

\subsection*{You'll want to know: \te{కవరు}}
\textbf{\te{కవరు}} — \emph{kavaru} — \textquotedblleft{}an envelope\textquotedblright{}.

\begin{grammarlens}[title={What you've built}]
One more word for this chapter.
\end{grammarlens}

\subsection*{Your turn}
\begin{itemize}
\raggedright
  \item \emph{Say it:} \emph{kavaru}
  \item \emph{Say it:} \emph{kavaru}, once more
  \item \emph{Say it:} say \emph{kavaru}
  \item \emph{Recall:} say the Telugu for a cucumber, then the Telugu for okra, then say \emph{kavaru} again
\end{itemize}

\begin{tcolorbox}[breakable,colback=teal!4,colframe=teal!35!black,title={Before you move on}]
What does \te{కవరు} mean? (\textquotedblleft{}An envelope\textquotedblright{}.) Say it once more.
\end{tcolorbox}
\section[\texorpdfstring{\te{స్టాంపు} (sṭāmpu)}{స్టాంపు (sṭāmpu)}]{\te{స్టాంపు} (sṭāmpu) — a postage stamp}

\label{lesson:TE-C185-stampu}



\begin{quote}
Before the new one: say the Telugu for okra, then the Telugu for an envelope.
\end{quote}

\subsection*{You'll want to know: \te{స్టాంపు}}
\textbf{\te{స్టాంపు}} — \emph{sṭāmpu} — \textquotedblleft{}a postage stamp\textquotedblright{}.

\begin{grammarlens}[title={What you've built}]
One more word for this chapter.
\end{grammarlens}

\subsection*{Your turn}
\begin{itemize}
\raggedright
  \item \emph{Say it:} \emph{sṭāmpu}
  \item \emph{Say it:} \emph{sṭāmpu}, once more
  \item \emph{Say it:} say \emph{sṭāmpu}
  \item \emph{Recall:} say the Telugu for okra, then the Telugu for an envelope, then say \emph{sṭāmpu} again
\end{itemize}

\begin{tcolorbox}[breakable,colback=teal!4,colframe=teal!35!black,title={Before you move on}]
What does \te{స్టాంపు} mean? (\textquotedblleft{}A postage stamp\textquotedblright{}.) Say it once more.
\end{tcolorbox}
\section[\texorpdfstring{\te{పార్సెలు} (pārselu)}{పార్సెలు (pārselu)}]{\te{పార్సెలు} (pārselu) — a parcel}

\label{lesson:TE-C185-parselu}



\begin{quote}
Before the new one: say the Telugu for an envelope, then the Telugu for a postage stamp.
\end{quote}

\subsection*{You'll want to know: \te{పార్సెలు}}
\textbf{\te{పార్సెలు}} — \emph{pārselu} — \textquotedblleft{}a parcel\textquotedblright{}.

\begin{grammarlens}[title={What you've built}]
One more word for this chapter.
\end{grammarlens}

\subsection*{Your turn}
\begin{itemize}
\raggedright
  \item \emph{Say it:} \emph{pārselu}
  \item \emph{Say it:} \emph{pārselu}, once more
  \item \emph{Say it:} say \emph{pārselu}
  \item \emph{Recall:} say the Telugu for an envelope, then the Telugu for a postage stamp, then say \emph{pārselu} again
\end{itemize}

\begin{tcolorbox}[breakable,colback=teal!4,colframe=teal!35!black,title={Before you move on}]
What does \te{పార్సెలు} mean? (\textquotedblleft{}A parcel\textquotedblright{}.) Say it once more.
\end{tcolorbox}
\section[\texorpdfstring{\te{తపాలా} (tapālā)}{తపాలా (tapālā)}]{\te{తపాలా} (tapālā) — the post, mail}

\label{lesson:TE-C185-tapala}



\begin{quote}
Before the new one: say the Telugu for a postage stamp, then the Telugu for a parcel.
\end{quote}

\subsection*{You'll want to know: \te{తపాలా}}
\textbf{\te{తపాలా}} — \emph{tapālā} — \textquotedblleft{}the post, mail\textquotedblright{}.

\begin{grammarlens}[title={What you've built}]
One more word for this chapter.
\end{grammarlens}

\subsection*{Your turn}
\begin{itemize}
\raggedright
  \item \emph{Say it:} \emph{tapālā}
  \item \emph{Say it:} \emph{tapālā}, once more
  \item \emph{Say it:} say \emph{tapālā}
  \item \emph{Recall:} say the Telugu for a postage stamp, then the Telugu for a parcel, then say \emph{tapālā} again
\end{itemize}

\begin{tcolorbox}[breakable,colback=teal!4,colframe=teal!35!black,title={Before you move on}]
What does \te{తపాలా} mean? (\textquotedblleft{}The post, mail\textquotedblright{}.) Say it once more.
\end{tcolorbox}
\section[\texorpdfstring{\te{పెన్సిలు} (pensilu)}{పెన్సిలు (pensilu)}]{\te{పెన్సిలు} (pensilu) — a pencil}

\label{lesson:TE-C185-pensilu}



\begin{quote}
Before the new one: say the Telugu for a parcel, then the Telugu for the post.
\end{quote}

\subsection*{You'll want to know: \te{పెన్సిలు}}
\textbf{\te{పెన్సిలు}} — \emph{pensilu} — \textquotedblleft{}a pencil\textquotedblright{}.

\begin{grammarlens}[title={What you've built}]
One more word for this chapter.
\end{grammarlens}

\subsection*{Your turn}
\begin{itemize}
\raggedright
  \item \emph{Say it:} \emph{pensilu}
  \item \emph{Say it:} \emph{pensilu}, once more
  \item \emph{Say it:} say \emph{pensilu}
  \item \emph{Recall:} say the Telugu for a parcel, then the Telugu for the post, then say \emph{pensilu} again
\end{itemize}

\begin{tcolorbox}[breakable,colback=teal!4,colframe=teal!35!black,title={Before you move on}]
What does \te{పెన్సిలు} mean? (\textquotedblleft{}A pencil\textquotedblright{}.) Say it once more.
\end{tcolorbox}
